% Part: incompleteness
% Chapter: introduction
% Section: definitions

\documentclass[../../../include/open-logic-section]{subfiles}

\begin{document}

\olfileid{inc}{int}{def}

\olsection{વ્યાખ્યાઓ}

ગણિતને ઔપચારિક બનાવવાનો અને આવું ઔપચારિકીકરણ
સુસંગત તથા પૂર્ણ છે એમ બતાવવાનો હિલ્બર્ટનો
કાર્યક્રમ હાથ ધરવા પહેલાં ભાષા, તાર્કિક માળખું
અને સ્વયંસિદ્ધ વાક્યોનું તંત્ર પસંદ કરવું પડે.
અહીં આપણે ધારી લઈએ કે ગણિતને પ્રથમ-ક્રમની
ભાષામાં ઔપચારિક બનાવી શકાય: એટલે, !!{constant}s,
!!{function}s અને !!{predicate}sનો કોઈ ગણ એવો
હોય કે પ્રથમ-ક્રમના તર્કશાસ્ત્રના સંયોજકો અને
પરિમાણકો સાથે મળીને તે ગણિતના દાવા વ્યક્ત
કરી શકે. મોટાભાગના લોકો માને છે કે એવી ભાષા
છે: ગણસિદ્ધાંતની ભાષા, જેમાં~$\in$ એકમાત્ર
અતાર્કિક ચિહ્ન છે. આટલી સાદી ભાષા આટલું બધું
વ્યક્ત કરી શકે છે એ પહેલી નજરે અવિશ્વસનીય
લાગે છે; ગણિતનો લગભગ આખો ભાગ આટલી
મર્યાદિત શબ્દાવલિમાં વ્યક્ત થઈ શકે છે તે
સ્થાપવામાં ઘણું કામ લાગ્યું. હમણાં માટે વાત
સરળ રાખવા અંકગણિત સુધી જ ચર્ચા મર્યાદિત
કરીએ, એટલે માત્ર પ્રાકૃતિક સંખ્યાઓ~$\Nat$
સાથે સંબંધિત ગણિત. અંકગણિતના તથ્યો વ્યક્ત
કરવાની સ્વાભાવિક ભાષા~$\Lang L_A$ છે. $\Lang L_A$માં
એક દ્વિસ્થાનિક !!{predicate}~$<$, એક
!!{constant}~$\Obj 0$, એક એકસ્થાનિક
!!{function}~$\prime$ અને બે દ્વિસ્થાનિક
!!{function}s~$+$ તથા~$\times$ છે.

\begin{defn}
!!{sentence}sનો ગણ~$\Gamma$ \emph{સિદ્ધાંત} છે, જો
તે તાર્કિક ફલિતતા હેઠળ બંધ હોય; એટલે કે જો
$\Gamma = \Setabs{!A}{\Gamma \Entails
!A}$ હોય.
\end{defn}

સિદ્ધાંતો નિર્દિષ્ટ કરવાની બે સરળ રીત છે. એક
રીત કોઈ !!{structure}માં સત્ય હોય એવાં
!!{sentence}sનો ગણ લેવાની. દાખલા તરીકે,
$\Lang L_A$ માટે એવું !!{structure} લો જેમાં
!!{domain}~$\Nat$ હોય અને બધાં અતાર્કિક
ચિહ્નોના અર્થઘટન અપેક્ષા મુજબ હોય.

\begin{defn}\ollabel{def:standard-model}
\emph{અંકગણિતનું પ્રમાણભૂત નિદર્શ} નીચે મુજબ
વ્યાખ્યાયિત !!{structure}~$\Struct N$ છે:
\begin{enumerate}
\item $\Domain N = \Nat$
\item $\Assign{\Obj 0}{N} = 0$
\item $\Assign{\Obj \prime}{N}(n) = n + 1$, દરેક $n \in \Nat$ માટે
\item $\Assign{\Obj +}{N}(n, m) = n + m$, દરેક $n, m \in \Nat$ માટે
\item $\Assign{\Obj \times}{N}(n, m) = n\cdot m$, દરેક $n, m \in \Nat$ માટે
\item $\Assign{\Obj <}{N} = \Setabs{\tuple{n, m}}{n \in \Nat, m \in
  \Nat, n < m}$
\end{enumerate}
\end{defn}

`$\times$' અને `$\cdot$'નો ભેદ નોંધો: $\times$
અંકગણિતની ભાષાનું ચિહ્ન છે. ગુણાકાર યાદ આવે
તે માટે એ પસંદ કર્યું છે, પણ $\times$ પોતે
ગુણાકારની ક્રિયા નથી; તે દ્વિસ્થાનિક વિધેય-ચિહ્ન
છે (સત્તાવાર રીતે~$\Obj f^2_1$). તેનાથી
વિપરીત~$\cdot$ સામાન્ય ગુણાકારનું વિધેય
\emph{છે}. $n \cdot m$ દેખાય ત્યારે તેનો
અર્થ $n$ અને~$m$ સંખ્યાઓનો ગુણાકાર છે;
$x \times y$ દેખાય ત્યારે અંકગણિતની ભાષાના
પદ વિશે વાત છે. પ્રમાણભૂત નિદર્શમાં વિધેય-ચિહ્ન
$\times$નું અર્થઘટન પ્રાકૃતિક સંખ્યાઓ પરના
વિધેય~$\cdot$ તરીકે થાય છે. સરવાળા માટે
આપણે~$+$ને અંકગણિતની ભાષાના વિધેય-ચિહ્ન
અને પ્રાકૃતિક સંખ્યાઓ પરના સરવાળાના વિધેય,
બંને માટે વાપરીએ છીએ. અહીં અર્થ પ્રસંગથી
નક્કી કરવો પડે.

\begin{defn}
\emph{સાચું અંકગણિત}નો સિદ્ધાંત અંકગણિતના
પ્રમાણભૂત નિદર્શમાં સંતોષાતાં !!{sentence}sનો ગણ છે;
એટલે કે,
\[
\Th{TA} = \Setabs{!A}{\Sat{N}{!A}}.
\]
\end{defn}

$\Th{TA}$ સિદ્ધાંત છે: $\Th{TA} \Entails !A$
હોય ત્યારે~$!A$, $\Th{TA}$ને સંતોષતા દરેક
!!{structure}માં સંતોષાય છે. $\Sat{N}{\Th{TA}}$
હોવાથી~$\Sat{N}{!A}$ મળે, તેથી
$!A \in \Th{TA}$.

સિદ્ધાંત~$\Gamma$ નિર્દિષ્ટ કરવાની બીજી રીત
તેને વાક્યોના કોઈ ગણ~$\Gamma_0$માંથી ફલિત
થતાં બધાં !!{sentence}sનો ગણ માનવાની છે.
એ સ્થિતિમાં~$\Gamma$, તાર્કિક ફલિતતા હેઠળ
$\Gamma_0$નું ``બંધ'' છે. આ રીતે સિદ્ધાંત
નિર્દિષ્ટ કરવો ત્યારે જ રસપ્રદ છે જ્યારે
$\Gamma_0$ સ્પષ્ટ રીતે આપવામાં આવ્યો હોય;
દાખલા તરીકે, $\Gamma_0$ના !!{element}sની યાદી હોય.
ઓછામાં ઓછું~$\Gamma_0$ નિર્ણેય તો હોવો
જોઈએ: કોઈ !!a{sentence} એ~$\Gamma_0$નો
ઘટક છે કે નહીં તેની સંગણકીય તપાસ શક્ય હોવી
જોઈએ. $\Gamma_0$નાં !!{sentence}sને આપણે
$\Gamma$નાં \emph{સ્વયંસિદ્ધ વાક્યો} અને
$\Gamma$ને~$\Gamma_0$ દ્વારા
\emph{સ્વયંસિદ્ધીકૃત} કહીએ છીએ.

\begin{defn}
સિદ્ધાંત~$\Gamma$ એ~$\Gamma_0$ વડે
\emph{સ્વયંસિદ્ધીકૃત} છે ત્યારે અને માત્ર ત્યારે,
જ્યારે
\[
\Gamma = \Setabs{!A}{\Gamma_0 \Entails !A}
\]
હોય.
\end{defn}

\begin{defn}
નીચેના વાક્યો વડે સ્વયંસિદ્ધીકૃત સિદ્ધાંત
$\Th{Q}$ ``રોબિન્સનનું~$\Th{Q}$'' કહેવાય છે;
તે અંકગણિતનો ખૂબ સરળ સિદ્ધાંત છે.
\begin{align*}
& \lforall[x][\lforall[y][(\eq[x'][y'] \lif \eq[x][y])]] \tag{$!Q_1$}\\
& \lforall[x][\eq/[\Obj 0][x']] \tag{$!Q_2$}\\
& \lforall[x][(\eq[x][\Obj 0] \lor \lexists[y][\eq[x][y']])] \tag{$!Q_3$}\\
& \lforall[x][\eq[(x + \Obj 0)][x]] \tag{$!Q_4$}\\
& \lforall[x][\lforall[y][\eq[(x + y')][(x + y)']]] \tag{$!Q_5$}\\
& \lforall[x][\eq[(x \times \Obj 0)][\Obj 0]] \tag{$!Q_6$}\\
& \lforall[x][\lforall[y][\eq[(x \times y')][((x \times y) + x)]]] \tag{$!Q_7$}\\
& \lforall[x][\lforall[y][(x < y \liff \lexists[z][\eq[(z' + x)][y]])]] \tag{$!Q_8$}
\end{align*}
!!{sentence}sનો ગણ $\{!Q_1, \dots, !Q_8\}$ એ
$\Th{Q}$નાં સ્વયંસિદ્ધ વાક્યો છે; એટલે આ
વાક્યોમાંથી ફલિત થતાં બધાં !!{sentence}s
$\Th{Q}$માં છે:
\[
\Th{Q} = \Setabs{!A}{\{!Q_1, \dots, !Q_8\} \Entails !A}.
\]
\end{defn}

\begin{defn}
માનો કે $!A(x)$ એ~$\Lang L_A$નું એવું
!!a{formula} છે જેના મુક્ત ચલ~$x$ અને
$y_1$, \dots, $y_n$ છે. તો નીચેના સ્વરૂપનું
કોઈ પણ !!{sentence}
\[
\lforall[y_1][\dots\lforall[y_n][((!A(\Obj 0) \land \lforall[x][(!A(x)
\lif !A(x'))]) \lif \lforall[x][!A(x)])]]
\]
\emph{આગમન યોજનાનું} એક ઉદાહરણ છે.

\emph{પિયાનો અંકગણિત}~$\Th{PA}$ એ
$\Th{Q}$નાં સ્વયંસિદ્ધ વાક્યો અને આગમન
યોજનાનાં બધાં ઉદાહરણોથી સ્વયંસિદ્ધીકૃત
સિદ્ધાંત છે.
\end{defn}

\begin{explain}
આગમન યોજનાનું દરેક ઉદાહરણ~$\Struct{N}$માં
સત્ય છે. !!{formula}~$!A$માં માત્ર એક મુક્ત
!!{variable}~$x$ હોય ત્યારે આ સૌથી સહેલું
દેખાય છે. ત્યારે~$!A(x)$, $\Struct{N}$માં
$\Nat$નો ઉપગણ~$X_{!A}$ વ્યાખ્યાયિત કરે છે.
$X_{!A}$ એ એવાં બધાં~$n \in \Nat$નો ગણ છે
જેમના માટે $s(x)=n$ હોય ત્યારે
$\Sat{N}{!A(x)}[s]$ છે. આગમન યોજનાનું
અનુરૂપ ઉદાહરણ આ છે:
\[
((!A(\Obj 0) \land \lforall[x][(!A(x) \lif !A(x'))]) \lif 
  \lforall[x][!A(x)]).
\]
તેનું પૂર્વાંગ~$\Struct{N}$માં સત્ય હોય,
તો $0 \in X_{!A}$ અને $n \in X_{!A}$ હોય
ત્યારે $n+1$ પણ તેમાં છે. $0 \in X_{!A}$
હોવાથી~$1 \in X_{!A}$ મળે; $1 \in X_{!A}$માંથી
$2 \in X_{!A}$ મળે; અને આમ આગળ.
તેથી દરેક~$n \in \Nat$ માટે
$n \in X_{!A}$. આનો અર્થ
$\lforall[x][!A(x)]$ એ~$\Struct{N}$માં
સંતોષાય છે.
\end{explain}

$\Th{Q}$ અને~$\Th{PA}$ બંને
સ્વયંસિદ્ધીકૃત સિદ્ધાંતો છે. મહત્ત્વનો પ્રશ્ન
એ છે કે તે કેટલા શક્તિશાળી છે? દાખલા તરીકે,
$\Th{PA}$, $\Lang L_A$માં વ્યક્ત થઈ શકે
એવાં~$\Nat$ વિશેનાં બધાં સત્યો સાબિત કરી
શકે છે? ચોક્કસ રીતે કહીએ, તો
$\Th{PA}$નાં સ્વયંસિદ્ધ વાક્યો
$\Lang L_A$માં પૂછાયેલા બધા પ્રશ્નો
નક્કી કરે છે?

બીજી રીતે પૂછીએ: શું $\Th{PA} = \Th{TA}$?
$\Th{TA}$ સ્પષ્ટપણે~$\Nat$ વિશેનાં બધાં
સત્યો સાબિત કરે છે (એટલે તેમને સમાવે છે) અને
$\Lang L_A$માં વ્યક્ત થઈ શકે એવા બધા
પ્રશ્નો નક્કી કરે છે. ખરેખર, $!A$ એ
$\Lang L_A$નું !!a{sentence} હોય તો કાં તો
$\Sat{N}{!A}$ અથવા~$\Sat{N}{\lnot !A}$;
એટલે કાં તો $\Th{TA} \Entails !A$ અથવા
$\Th{TA} \Entails \lnot !A$. આવા સિદ્ધાંતને
\emph{!!{complete}} કહીએ.

\begin{defn}
સિદ્ધાંત~$\Gamma$ \emph{!!{complete}} છે
ત્યારે અને માત્ર ત્યારે, જ્યારે તેની ભાષાના
દરેક !!{sentence}~$!A$ માટે કાં તો
$\Gamma \Entails !A$ અથવા
$\Gamma \Entails \lnot !A$ હોય.
\end{defn}

\begin{explain}
પૂર્ણતાના પ્રમેય મુજબ
$\Gamma \Entails !A$ ત્યારે અને માત્ર ત્યારે
જ્યારે $\Gamma \Proves !A$. તેથી~$\Gamma$
પૂર્ણ છે ત્યારે અને માત્ર ત્યારે, જ્યારે
તેની ભાષાના દરેક !!{sentence}~$!A$ માટે
કાં તો $\Gamma \Proves !A$ અથવા
$\Gamma \Proves \lnot !A$ હોય.
\end{explain}

બીજો પ્રશ્ન પણ ઊભો થાય છે: કોઈ
!!a{sentence} એ~$\Th{TA}$, $\Th{PA}$
અથવા માત્ર~$\Th{Q}$માં પણ છે કે નહીં તે
તપાસવાની સંગણકીય પ્રક્રિયા છે? કોઈ ગણ
(દાખલા તરીકે, !!{sentence}sનો ગણ)
ક્યારે !!{decidable} છે તેની વ્યાખ્યા
આપીને પ્રશ્ન વધુ ચોક્કસ કરી શકીએ.

\begin{defn}
ગણ~$X$ \emph{!!{decidable}} છે ત્યારે
અને માત્ર ત્યારે, જ્યારે એવી સંગણકીય પ્રક્રિયા
હોય જે ઇનપુટ~$x$ માટે $x \in X$ હોય તો
$1$ અને અન્યથા~$0$ આપે.
\end{defn}

એટલે પ્રશ્ન આ છે: $\Th{TA}$
($\Th{PA}$, $\Th{Q}$) !!{decidable} છે?

આ બધા પ્રશ્નોનો જવાબ ``ના'' છે. આમાંથી
કોઈ સિદ્ધાંત નિર્ણેય નથી. જોકે આ બાબત માત્ર
આ ખાસ સિદ્ધાંતો પૂરતી નથી. ખરેખર,
ચોક્કસ શરતો સંતોષતા \emph{દરેક} સિદ્ધાંતને
એવાં જ પરિણામો લાગુ પડે છે. $\Th{Q}$
અને~$\Th{PA}$ને લાગુ પડતી એવી એક શરત
એ છે કે તેઓ સ્વયંસિદ્ધ વાક્યોના
!!a{decidable} ગણથી સ્વયંસિદ્ધીકૃત છે.

\begin{defn}
સિદ્ધાંત \emph{!!{axiomatizable}} છે, જો
તે સ્વયંસિદ્ધ વાક્યોના !!a{decidable} ગણથી
સ્વયંસિદ્ધીકૃત હોય.
\end{defn}

\begin{ex}
!!{sentence}sના સાન્ત ગણથી સ્વયંસિદ્ધીકૃત
દરેક સિદ્ધાંત !!{axiomatizable} છે, કારણ કે
દરેક સાન્ત ગણ !!{decidable} છે. તેથી,
દાખલા તરીકે, $\Th{Q}$ !!{axiomatizable} છે.

$\Th{PA}$ જેવા યોજના વડે સ્વયંસિદ્ધીકૃત
સિદ્ધાંતો પણ !!{axiomatizable} છે. $!B$ એ
$\Th{PA}$ના સ્વયંસિદ્ધ વાક્યોમાંનું છે કે
નહીં તે તપાસવા—એટલે કે $!B$ એ
$\Th{PA}$નું સ્વયંસિદ્ધ વાક્ય હોય તો
$\Char{X}(!B)=1$ અને અન્યથા~$=0$
આપતું વિધેય~$\Char{X}$ ગણવા—આ કરી શકાય:
પહેલાં તપાસો કે $!B$ એ~$\Th{Q}$ના
સ્વયંસિદ્ધ વાક્યોમાંનું એક છે કે નહીં.
હોય તો જવાબ ``હા'' અને
$\Char{X}(!B)=1$. ન હોય તો તે આગમન
યોજનાનું ઉદાહરણ છે કે નહીં તપાસો. આ
તપાસ પદ્ધતિસર થઈ શકે છે. અહીં તે આ રીતે
થાય છે એમ જોવું કદાચ સહેલું છે: આગમન
યોજનાનું કોઈ પણ ઉદાહરણ થોડા સર્વપરિમાણકોથી
શરૂ થાય છે અને પછી શરતી ઉપ-!!{formula}
આવે છે. તે શરતનું ઉત્તરાંગ
$\lforall[x][!A(x, y_1, \dots, y_n)]$
છે, જ્યાં $x$ અને $y_1$, \dots, $y_n$
એ બધાં~$!A$ના મુક્ત ચલ છે; $!B$ના
પ્રારંભિક પરિમાણકો $y_1$, \dots,~$y_n$
ને બાંધે છે. આ~$!A$ કાઢીને તેના મુક્ત
ચલો પ્રારંભના સર્વપરિમાણકો અને
$\lforall[x]$થી બંધાયેલા ચલો સાથે સરખા
છે એ તપાસ્યા પછી શરતનું પૂર્વાંગ
નીચેના સૂત્ર સાથે મેળ ખાય છે કે નહીં
તે તપાસીએ:
\[
!A(\Obj 0, y_1, \dots, y_n) \land \lforall[x][(!A(x, y_1, \dots, y_n)
\lif !A(x', y_1, \dots, y_n))]
\]
મેળ ખાય તો~$!B$ આગમન યોજનાનું
ઉદાહરણ છે; ન ખાય તો~$!B$ ઉદાહરણ નથી.
\end{ex}

આ પ્રશ્નનો—અને કયા સિદ્ધાંતો પૂર્ણ અથવા
નિર્ણેય છે એના વધુ સામાન્ય પ્રશ્નનો—જવાબ
શોધવામાં નીચેની વ્યાખ્યા પણ ઉપયોગી થશે.
યાદ કરો કે ગણ~$X$ !!{enumerable} છે ત્યારે
અને માત્ર ત્યારે, જ્યારે તે ખાલી હોય અથવા
કોઈ !!a{surjective} વિધેય
$f \colon \Nat \to X$ હોય. આવું વિધેય
$X$ની પરિગણના કહેવાય છે.

\begin{defn}
ગણ~$X$ \emph{!!{computably enumerable}}
(ટૂંકમાં~!!{c.e.}) કહેવાય છે ત્યારે અને માત્ર
ત્યારે, જ્યારે તે ખાલી હોય અથવા તેની
સંગણનીય પરિગણના હોય.
\end{defn}

!!{axiomatizability} ઉપરાંત અપૂર્ણતાનાં
પ્રમેયો લાગુ પડે એવા સિદ્ધાંતો માટે બીજી શરત
એ છે કે તેઓ સંગણનીય વિધેયો અને
!!{decidable} સંબંધો વિશેનાં મૂળભૂત તથ્યો
સાબિત કરવા જેટલા શક્તિશાળી હોય. ``મૂળભૂત
તથ્યો'' એટલે વિધેયના દરેક નિવેશ માટે તેનું
મૂલ્ય વ્યક્ત કરતાં !!{sentence}s. અને
``પૂરતા શક્તિશાળી'' એટલે આ વાક્યો એ
સિદ્ધાંતોના પ્રમેયોમાં હોય. દાખલા તરીકે
સરવાળો એક સામાન્ય સંગણનીય વિધેય છે.
$2$ અને~$3$ નિવેશ માટે~$+$નું મૂલ્ય~$5$
છે, એટલે $2+3=5$. અંકગણિતની ભાષામાં
$2$ અને~$3$ નિવેશ માટે~$+$નું મૂલ્ય~$5$
છે તે વ્યક્ત કરતું વાક્ય છે:
$(\num{2} + \num{3}) = \num{5}$.
દાખલા તરીકે, $\Th{Q}$ આ વાક્ય સાબિત
કરે છે. વધુ સામાન્ય રીતે, દરેક સંગણનીય
વિધેય~$f(x_1, x_2)$ માટે~$\Lang{L_A}$માં
એવું !!a{formula}~$!A_f(x_1, x_2, y)$
હોવું જોઈએ કે જ્યારે $f(n_1, n_2)=m$
હોય ત્યારે $\Th{Q} \Proves
!A_f(\num{n_1}, \num{n_2}, \num{m})$.
આ રીતે $\Th{Q}$ સાબિત કરે છે કે
$n_1$, $n_2$ નિવેશ માટે~$f$નું મૂલ્ય~$m$
છે. વાસ્તવમાં, આપણે થોડું વધુ માગીએ છીએ:
$n_1$, $n_2$ નિવેશ માટે આ મૂલ્ય સિવાય
બીજી કોઈ સંખ્યા~$f$નું મૂલ્ય નથી, તે પણ
સાબિત થાય.
!!{decidable} સંબંધો માટે પણ આવી જ શરત
છે. આગળની બે વ્યાખ્યાઓમાં આ ચોક્કસ કરીએ.

\begin{defn}
!!{formula}~$!A(x_1, \dots, x_k, y)$ એ
વિધેય~$f\colon \Nat^k \to \Nat$ને
$\Gamma$માં \emph{!!{represents}} છે
ત્યારે અને માત્ર ત્યારે, જ્યારે
$f(n_1, \dots, n_k)=m$ હોય ત્યારે
\begin{enumerate}
\item $\Gamma \Proves !A(\num{n_1}, \dots, \num{n_k}, \num{m})$, અને
\item $\Gamma \Proves \lforall[y](!A(\num{n_1}, \dots, \num{n_k},
y) \lif y = \num{m})$.
\end{enumerate}
\end{defn}

\begin{defn}
!!{formula}~$!A(x_1, \dots, x_k)$ સંબંધ
$R \subseteq \Nat^k$ને એ જ સિદ્ધાંતમાં
\emph{!!{represents}}
છે ત્યારે અને માત્ર ત્યારે,
\begin{enumerate}
\item $R(n_1, \dots, n_k)$ હોય ત્યારે
  $\Gamma \Proves !A(\num{n_1},
\dots, \num{n_k})$, અને
\item $R(n_1, \dots, n_k)$ ન હોય ત્યારે
  $\Gamma \Proves \lnot
!A(\num{n_1}, \dots, \num{n_k})$.
\end{enumerate}
\sourcecorrection{OLINC-002}{સ્થિર અંગ્રેજી મૂળમાં આ
સંબંધના નિરૂપણની વ્યાખ્યામાં સિદ્ધાંત સ્પષ્ટ કહ્યો
નથી, છતાં બંને શરતોમાં~$\Gamma$ વપરાય છે.
અહીં ``એ જ સિદ્ધાંતમાં''થી અગાઉની વ્યાખ્યાનો
સિદ્ધાંત ઉદ્દિષ્ટ કર્યો છે; સૂત્રો બદલ્યાં નથી.}
\end{defn}

સિદ્ધાંત અપૂર્ણતાનાં પ્રમેયો લાગુ પડે એટલો
``પૂરતો શક્તિશાળી'' છે, જો તે બધાં સંગણનીય
વિધેયો અને બધાં !!{decidable} સંબંધોને
!!{represents} કરે. $\Th{Q}$ અને તેના
વિસ્તારો આ શરત સંતોષે છે, પરંતુ આ સ્થાપવામાં
આપણને સમય લાગશે: $\Th{Q}$ કેવાં તથ્યો
સાબિત કરી શકે તે વિશે આ અ-સરળ પરિણામ
છે, અને $\Th{Q}$નાં થોડાં જ સ્વયંસિદ્ધ વાક્યોમાંથી
બધું સાબિત કરવું પડે છે. તેમ છતાં
$\Th{Q}$ ખૂબ નબળો સિદ્ધાંત છે. એટલે
$\Th{Q}$ બધાં સંગણનીય વિધેયોનું
નિરૂપણ કરે છે તે સાબિત કરવું મુશ્કેલ
હોવા છતાં, મોટાભાગના રસપ્રદ સિદ્ધાંતો
$\Th{Q}$ કરતાં વધુ શક્તિશાળી છે, એટલે
$\Th{Q}$ કરતાં વધુ સાબિત કરે છે.
$\Th{Q}$ કંઈક સાબિત કરે તો દરેક વધુ
શક્તિશાળી સિદ્ધાંત પણ કરે. $\Th{Q}$
બધાં સંગણનીય વિધેયોનું નિરૂપણ કરે છે,
તેથી દરેક વધુ શક્તિશાળી સિદ્ધાંત પણ
કરે છે. આમ, ઘણા રસપ્રદ સિદ્ધાંતો
અપૂર્ણતાનાં પ્રમેયોની આ શરત સંતોષે છે.
એટલે આપણું મુશ્કેલ કામ ઉપયોગી નીવડશે:
તે આ પ્રમેયો ઘણા સિદ્ધાંતોને લાગુ પડે
છે તે બતાવે છે. ``સમગ્ર ગણિત''ને
ઔપચારિક બનાવવાનો હેતુ ધરાવતો કોઈ
સિદ્ધાંત $\Th{Q}$ જે સાબિત કરે છે
તે બધું તો સાબિત કરી શકે જ જોઈએ;
ઓછામાં ઓછું પ્રાથમિક સંગણનાઓના
પરિણામોને તે સમાવી શકે. તેથી
``સમગ્ર ગણિત''નો ઉમેદવાર સિદ્ધાંત
અપૂર્ણતાનાં પ્રમેયો લાગુ પડે એવો હશે.


\end{document}
